\SourceLeaf{303}%
Punkte $A$, $B$, $C$, $D$ (ich brauche jetzt meine, nicht \textsc{Newtons} hierbei
meines Erachtens unbehülfliche Bezeichnung) einen Punkt $\delta$ jener geraden Linie
sucht, dann vermittelst $A$, $B$, $C$, $E$ einen andern Punkt $\varepsilon$ jener
geraden Linie; so ist die gerade Linie selbst positione data, und für jeden andern
Punkt derselben, $\mu$, gibt seine Construction einen correspondirenden
Kegelschnittspunkt $M$: indem er den Winkel $BAC$ (mit unbestimmter Schenkellänge)
sich um $A$ drehen lässt, und die Stellung notirt, wo $AB$ durch $\mu$ geht, muss $AC$
durch $M$ gehen, und eben so, wenn der Winkel $ABC$ sich um $B$ dreht, muss, indem
$BA$ durch $\mu$ geht, $BC$ gleichfalls durch $M$ gehen. Diese Construction ist nun
zwar überaus zierlich, aber Herr des Gegenstandes ist man doch erst dann, wenn man
alle andern diese magische gerade Linie betreffenden Fragen beantworten kann;
namentlich will man wissen, welche Relationen diese gerade Linie zu den Elementen des
Kegelschnitts habe, ob man diese Elemente selbst mit Leichtigkeit aus der Lage jener
geraden Linie und der von $A$, $B$, $C$ ableiten könne. Verschiedenes dieser Art kann
ich jetzt recht artig ausrichten, ich weiss aber nicht, ob ich selbst das Ganze
durchführen kann, da andere Geschäfte mich nöthigen abzubrechen.

\mbox{.\,.\,.\,.\,.\,.\,.} Früher haben Sie, wenn ich nicht irre, sich viel mit den Kegelschnitten
beschäftigt; gewiss viel mehr darüber gelesen als ich. Sie führen z. B. schon
Monatliche Correspondenz 1810 November S. 508 einen eleganten die Ellipse betreffenden
Satz aus einem Buche\HistoricalNote{[*)]}{[*)}{\textsc{Puissant}, Recueil de diverses propositions de géométrie. Paris 1801.
Vergl. auch Bd. IV. S. 396.]} an, welches ich, obwohl vermuthlich unsere Bibliothek es
besitzt, bis jetzt noch nicht angesehen habe. Man kann also wegen historischer Notiz
wohl bei Ihnen vorfragen. Kennen Sie eine zierliche Construction, um den \textls{Mittelpunkt}
eines Kegelschnitts aus 5 Punkten desselben --- \textls{ohne alle überflüssigen Umwege} --- zu
finden? Schwer wird es gewiss nicht sein, aber gerade diese Entwickelung habe ich noch
nicht \mbox{gemacht. .\,.\,.\,.\,.\,.\,.}
